\chapter{Discussion}
\label{ch:discussion}

\section{When the Method Does Not Help}
\label{sec:disc-limits}

A method is more useful when its failure conditions are written down than when
the next person discovers them. The scheduler beats the incumbent almost
everywhere, but the part of it that is \emph{new} --- the division by cost ---
needs three conditions, and the pilot gives evidence on all three.

\begin{table}[htbp]
\centering
\caption{The three conditions under which the cost term is expected to pay,
and what the method degrades to when each fails.}
\label{tab:conditions}
\small
\begin{tabularx}{\textwidth}{@{}>{\raggedright\arraybackslash}p{0.21\textwidth} X X@{}}
\toprule
Condition & Measured by & What happens when it fails\\
\midrule
\textbf{C1. The pool is large enough} &
$|P|$, known before any run &
With a handful of tests there is little to order, and a candidate that survives
pays for the whole pool whatever the order. Every configuration converges to
one schedule.\\
\addlinespace
\textbf{C2. Test costs are not uniform} &
The within-task coefficient of variation of per-test CPU
(Table~\ref{tab:res-arena}) &
$\hat{m}$ becomes a constant, dividing by it changes no ranking, and the rule
collapses to \frtp\ with a denominator.\\
\addlinespace
\textbf{C3. The task is hard enough for the proposer} &
How many rounds the loop needs on that task &
A task repaired in one or two rounds never builds a history, so there is
nothing to schedule by and nothing to select hints from.\\
\bottomrule
\end{tabularx}
\end{table}

\paragraph{C2 is the one the pilot can speak to, and it holds up.} The
within-task cost spread is the variable that splits
Table~\ref{tab:res-cv}, and the comparison against \frtp\ runs
$+1.2\,\%$, $+1.1\,\%$, $-4.1\,\%$ across the three groups. That is the
predicted ordering: the cost term costs about one percent where there is no
spread and is the only group where it wins where there is. None of the three
intervals excludes zero, so this is the \emph{shape} of the prediction
confirmed rather than the prediction itself.

The quantity that matters is the spread \emph{within} a task, because a
schedule only ever chooses among the tests of one task: a corpus of
individually uniform tasks looks heterogeneous in aggregate and offers the
scheduler nothing. Here the aggregate spread is enormous --- the
99th-percentile test costs $177$ times the median --- while the median
within-task CV is only $0.159$, so two thirds of these tasks are in the regime
where C2 fails.

\paragraph{C3 dilutes rather than harms.} A corpus of easy tasks does not make
the method look bad; it makes every configuration look identical and pulls a
corpus-wide average towards no effect --- visible here in the $34$ of $97$
tasks on which \code{sched} and \code{frtp} cost exactly the same. Difficulty
is therefore a stratification variable, not a caveat.

Note that C1 and C2 are properties of the frozen pool, known before any
candidate exists, so splitting on them stratifies a fixed population. C3 is an
\emph{outcome}, and admitting tasks on the basis of it would select the corpus
on the result --- so the pilot reports the whole-corpus number as the headline
and the split as a secondary reading, never the other way round.

\section{Reporting Seconds Rather Than Executions}
\label{sec:disc-units}

The prioritisation literature reports the number of test executions rather than
elapsed time, for a good reason: the largest study of patch validation to date
counts executions precisely because elapsed time is unstable across
environments~\citep{lou2024aprregression}. Time depends on the machine, on what
else is running, and on the interpreter; a count does not.

The problem is that the reproducible unit and the deployed unit are different
quantities, and the choice between them is not neutral when test costs vary.
The pilot produced the textbook case. Against \frtp, the proposed scheduler
runs $13\,\%$ \emph{more} test executions and spends $3\,\%$ \emph{less} CPU.
Reported in executions it lost; reported in seconds it won. A deployment pays
for seconds.

This report therefore gives \textbf{both units for every configuration} --- it
costs nothing, since a loop that measures time is already counting executions
--- and reports CPU rather than wall clock, which on a machine running several
workers is contaminated by siblings. It also reports the \textbf{ratio}, CPU
per execution, because that is what separates ``found the rejection in fewer
steps'' from ``found it in cheaper steps''. In Table~\ref{tab:res-rq1b} that
column is the whole mechanism: $0.401$\,s per run for \frtp\ against $0.344$\,s
for \code{sched}.

None of this contradicts \citeauthor{lou2024aprregression}'s reasoning. Their
setting --- Java projects with method-level suites --- is one where test
\emph{selection} is available: a test can be skipped entirely because no
changed method reaches it, and there they find selection to dominate
prioritisation. Here a test is one whole-program execution on one input, so
there is no call graph to cut and nothing to select away. Prioritisation is the
only lever.

\section{Composing with the Other Two Families, and What Each Costs to Adopt}
\label{sec:disc-combining}

The method changes only two things: which test runs first, and what goes into
the prompt. It does not change how a test is executed, and it does not change
whether a candidate is executed at all. That makes it composable with the other
two families of validation-cost reduction, each of which reduces a different
factor of Equation~\eqref{eq:three-factors}.

\textbf{Execution acceleration} makes each run cheaper ---
\textsc{UniAPR}~\citep{chen2021uniapr} and
\textsc{ExpressAPR}~\citep{xiao2024expressapr} reuse a warm runtime and adapt
mutation-testing techniques, reporting up to $137\times$. A scheduler is
indifferent to how a test is executed: it ranks by measured cost, so if
acceleration changes the costs it re-learns them. One interaction cuts in a
useful direction: acceleration typically removes a fixed per-run overhead,
which \emph{increases} the relative spread between cheap and expensive tests,
and that spread is exactly what the cost term needs. \textbf{Execution-free
verification} removes candidates before they are run;
\textsc{R2E-Gym}~\citep{jain2025r2egym} finds learned and test-based verifiers
complementary, so a learned verifier can filter cheaply and the scheduler then
decides how to spend execution on the survivors. \textbf{Selected hints} also
reduce the number of candidates, but earlier: they change what is generated
rather than filtering it afterwards.

\paragraph{They do not cost the same to adopt.} Acceleration needs an
integration with one language runtime --- class reloading, state resetting ---
and the correctness argument that comes with reusing a dirty runtime.
Execution-free verification needs training data, a training run, a second model
to serve, and a standing question about distribution shift. Search over repair
attempts~\citep{antoniades2025swesearch} spends extra model calls per bug,
which increases the resource this work is trying to reduce. The scheduler needs
three counters, one sort per round, and a timer around a call the loop already
makes (Section~\ref{sec:method-cost}).

It is therefore the cheapest of the four to try and the easiest to withdraw:
if it does not help it degrades to the ordering the loop already used, and
nothing has to be retrained when the proposer is swapped, because the
statistics are of the task and the session rather than of the model. That is
also an argument about where to start. Acceleration and verifiers have larger
ceilings and both require a project; a scheduler is an afternoon, and it
establishes whether the pool has any cost structure worth exploiting before
anyone commits to the larger work --- $60\,\%$ of validation CPU here was
removed by three counters.

\paragraph{Simplicity is a claim that has to survive the ablation, and it only
partly did.} Table~\ref{tab:res-ladder} shows a plain rejection count ---
\frtp, published in $2013$ --- already captures nearly the whole effect here.
The honest statement of the contribution is therefore narrower than the
headline: scheduling matters enormously, and the cost term adds a further
amount that grows over a session and concentrates where the theory says it
should. Establishing that the second part is real needs the full corpus.

\section{What the Pilot Can and Cannot Settle}
\label{sec:disc-scope}

A pilot of $97$ tasks is a screening instrument: it can show that an effect
exists and is large enough to justify a corpus run, and fix the parameters
that run must declare in advance. It cannot establish a magnitude.

\textbf{Settled:} scheduling the oracle is worth doing ($-60.8\,\%$, no cost
in repairs), and memory rather than the score is the dominant term.
\textbf{Not settled:} whether dividing by cost beats \frtp\ at corpus scale,
whether a rate beats a count at a larger budget, and anything on the
generation side.

One caution applies to the cost exponent $c$ and the hint count $k$: a value
selected on the data it is reported on is optimistically biased, so both are
declared before the corpus run (Section~\ref{sec:future}).
